\documentclass[5pt]{article}
\usepackage{multicol,multirow}
\usepackage{graphicx} % Required for inserting images
\usepackage[margin=0.45cm]{geometry}
\usepackage{xcolor}
\usepackage{array}

\definecolor{LightGray}{gray}{0.9}
\usepackage[inline]{enumitem}
\usepackage{minted}

\newcolumntype{L}{>{\footnotesize}l}
\newcolumntype{C}{>{\ttfamily\footnotesize}l}

\begin{document}
\begin{center}
    \Large{\textbf{SQLite Cheat Sheet}}\\
    \small{Last Updated: \today}\hfill\small{\textcopyright Maximilien Notz \the\year{}}
    \hspace*{-\dimexpr\oddsidemargin+1in\relax}\rule{\paperwidth}{0.4pt}
\end{center}

% content 
\begin{multicols}{2}
\setcounter{secnumdepth}{0}

\subsection{Clauses}
\textbf{Clause order:} \texttt{SELECT} $\rightarrow$ \texttt{FROM} $\rightarrow$ \texttt{WHERE} $\rightarrow$ \texttt{GROUP BY} $\rightarrow$ \texttt{HAVING} $\rightarrow$ \texttt{ORDER BY} $\rightarrow$ \texttt{LIMIT}\\
\begin{tabular}{C L}
    SELECT x          & Selects the column \texttt{x} from the table.\\
    SELECT *          & Selects all columns from the table.\\
    SELECT DISTINCT x & Selects unique values from the column \texttt{x}.\\
    FROM table        & Specifies the table to select from.\\
    WHERE x           & Filters the results based on the condition \texttt{x}.\\
    ORDER BY x        & Sorts the results based on the column \texttt{x}.\\
    GROUP BY x        & Groups the results based on the column \texttt{x}.\\
    AS x              & Renames a column or table to \texttt{x}.\\
    HAVING x          & Filters the results of a \texttt{GROUP BY} clause based on the\\
                      & condition \texttt{x}.\\
    LIKE x            & Used in \texttt{WHERE} clauses to match a pattern.\\
                      & "\%" many characters, "\_" single character\\
    LIMIT x           & Limits the number of rows returned to \texttt{x}.\\
\end{tabular}

\subsection{Table Modifiers}
It is a good practice to Select before modifying a table.\\
\begin{tabular}{C L}
    INSERT INTO x(cols) VALUES (lst) & Make a row-wise insertion of the\\
                                     & values in \texttt{lst} into the table \texttt{x}.\\
    CREATE TABLE x (cols)            & Creates a new table named \texttt{x} with the specified columns.\\
    DELETE FROM x WHERE c            & Deletes rows from the table \texttt{x} where\\
                                     & the \texttt{c} is met.\\
    DROP TABLE x                     & Deletes the table \texttt{x}.\\
    UPDATE x SET c=v WHERE c         & Updates the column \texttt{c} in table \texttt{x} to\\
                                     & the value \texttt{v} if the \texttt{c} is met.\\
\end{tabular}
In create queries we can specify the data type and constraints of each column. The most common data types are:
\begin{tabular}{C L}
    INTEGER & A signed integer.\\
    REAL    & A floating point value.\\
    TEXT    & A text string.\\
    BLOB    & A binary large object.\\
\end{tabular}\\
And the most common constraints are:\\
\begin{tabular}{C L}
    PRIMARY KEY & Uniquely identifies each row in the table.\\
    FOREIGN KEY & Establishes a relationship between two tables.\\
    NOT NULL    & Ensures that a column cannot have a \texttt{NULL} value.\\
    UNIQUE      & Ensures that all values in a column are unique.\\
    CHECK       & Ensures that all values in a column satisfy a specific condition.\\
    DEFAULT     & Sets a default value for a column when no value is specified.\\
\end{tabular}

\subsection{Operator}
\begin{tabular}{C L}
    LIKE x          & Used in \texttt{WHERE} clauses to match a pattern.\\
                    & \texttt{\%} represents zero, one, or multiple characters.\\
                    & \texttt{\_} represents a single character.\\
    AND x           & Combines multiple conditions in a \texttt{WHERE} clause.\\
    OR x            & Combines multiple conditions in a \texttt{WHERE} clause.\\
    NOT x           & Negates a condition in a \texttt{WHERE} clause.\\
    IN (x, y)       & Checks if a value is in a list of values.\\
    BETWEEN x AND y & Checks if a value is within a range of values.\\
    IS NULL         & Checks if a value is \texttt{NULL}.\\
    IS NOT NULL     & Checks if a value is not \texttt{NULL}.\\
\end{tabular}

\subsection{Functions}
\begin{tabular}{C L}
    TOP(n)      & Limits the number of rows returned to \texttt{n}.\\
    COUNT(x)    & Counts the number of non-null values in column \texttt{x}.\\
    SUM(x)      & Calculates the sum of values in column \texttt{x}.\\
    AVG(x)      & Calculates the average of values in column \texttt{x}.\\
    MIN(x)      & Finds the minimum value in column \texttt{x}.\\
    MAX(x)      & Finds the maximum value in column \texttt{x}.\\
\end{tabular}

\subsection{Examples}
\begin{minted}[bgcolor=LightGray]
{sql}
SELECT Orders.OrderID, Customers. FirstName,
Customers.LastName, Customers.Email, Orders.OrderDate
FROM Orders
JOIN Customers ON Orders.CustomerID=Customers.CustomerID
ORDER BY OrderID
\end{minted}

\subsubsection{Nested Queries}
\begin{minted}[bgcolor=LightGray]
{sql}
SELECT * FROM Orders
WHERE CustomerID=(
	SELECT CustomerID FROM Customers
	WHERE Email="ex@example.com"
)
\end{minted}
\subsubsection{Insert Queries}
\begin{minted}[bgcolor=LightGray]{sql}
INSERT INTO EMPLOYEES(NAME, FIRSTNAME, DEPT,DEPTFLOOR) 
VALUES ('Brown', 'Tom', 'Management', 1)
\end{minted}

\subsubsection{Table Creation}
\begin{minted}[bgcolor=LightGray]{sql}
CREATE TABLE EMPLOYEES(
    "id" INTEGER NOT NULL,
    "name" TEXT NOT NULL,
    "first_name" TEXT NOT NULL,
    "dept_id" TEXT NOT NULL,
    "dept_floor" INTEGER NOT NULL,
    PRIMARY KEY("id"),
    FOREIGN KEY("dept") 
      REFERENCES "Departments"("dept_id")
)
\end{minted}


\end{multicols}
\end{document}
